\section*{Problemas}

\subsection*{Enunciados de los problemas}

% Recordar
\begin{problema}
	\label{prob:20ee012}
	Enuncie, sin realizar ningún cálculo, la definición de $F=(\chi^2_{d_1}/d_1)/(\chi^2_{d_2}/d_2)$ con $\chi^2_{d_1}$ y $\chi^2_{d_2}$ independientes, y la fórmula de $E(F)$.
\end{problema}

% Comprender
\begin{problema}
	\label{prob:2e8078c}
	Sea $F \sim F_{5,8}$. Explique, sin calcular ninguna probabilidad, por qué $1/F$ sigue una distribución $F_{8,5}$ (es decir, por qué invertir el cociente simplemente intercambia los grados de libertad del numerador y del denominador).
\end{problema}

% Aplicar
\begin{problema}
	\label{prob:b47df2a}
	Dos líneas de producción se comparan: la línea 1 ($n_1=13$) tiene $S_1^2=45$, y la línea 2 ($n_2=11$) tiene $S_2^2=20$. Pruebe $H_0:\sigma_1^2=\sigma_2^2$ a nivel $\alpha=0.05$ (bilateral), usando $F_{12,10,0.025}\approx3.62$.
\end{problema}

% Analizar
\begin{problema}
	\label{prob:23d903e}
	Sea $T \sim t_\nu$. Demuestre que $T^2 \sim F_{1,\nu}$, partiendo de la definición $T=Z/\sqrt{\chi^2_\nu/\nu}$ con $Z\sim N(0,1)$ y $Z$ independiente de $\chi^2_\nu$.
\end{problema}

% Evaluar
\begin{problema}
	\label{prob:e04a795}
	Dos muestras independientes tienen $n_1=15$, $S_1^2=50$ y $n_2=12$, $S_2^2=48$. Primero pruebe $H_0:\sigma_1^2=\sigma_2^2$ con la prueba $F$ (use $F_{14,11,0.025}\approx3.36$). Con base en el resultado, evalúe si sería apropiado usar la prueba $t$ agrupada (\emph{pooled}) o la prueba $t$ de Welch (varianzas desiguales) para comparar las medias, y justifique.
\end{problema}

% Crear
\begin{problema}
	\label{prob:6277cc2}
	Diseñe un ejemplo original (contexto, dos muestras con sus tamaños y varianzas) donde se aplique la prueba $F$ para comparar dos varianzas poblacionales, distinto de los ejemplos de líneas de producción ya vistos en esta sección. Calcule el estadístico $F$ y establezca la decisión usando un valor crítico razonable.
\end{problema}

\subsection*{Soluciones detalladas}

% Recordar
\begin{solproblema}[prob:20ee012]
	$F = \dfrac{\chi^2_{d_1}/d_1}{\chi^2_{d_2}/d_2}$, con $\chi^2_{d_1}$ y $\chi^2_{d_2}$ independientes. Su media es $E(F) = \dfrac{d_2}{d_2-2}$ (para $d_2>2$).
\end{solproblema}

% Comprender
\begin{solproblema}[prob:2e8078c]
	Si $F = (\chi^2_{d_1}/d_1)/(\chi^2_{d_2}/d_2)$, entonces $1/F = (\chi^2_{d_2}/d_2)/(\chi^2_{d_1}/d_1)$: el mismo cociente, pero con el papel de "numerador" y "denominador" intercambiado entre las dos chi-cuadradas. Como la definición de $F_{a,b}$ es precisamente "chi-cuadrada con $a$ grados de libertad (normalizada) sobre chi-cuadrada con $b$ grados de libertad (normalizada)", invertir $F_{d_1,d_2}$ produce exactamente esa misma estructura con los grados de libertad intercambiados, es decir, $F_{d_2,d_1}$.
\end{solproblema}

% Aplicar
\begin{solproblema}[prob:b47df2a]
	$F = 45/20 = 2.25$. Como $2.25 < 3.62 = F_{12,10,0.025}$, \textbf{no se rechaza} $H_0:\sigma_1^2=\sigma_2^2$: no hay evidencia suficiente de que las varianzas difieran.
\end{solproblema}

% Analizar
\begin{solproblema}[prob:23d903e]
	$T^2 = \dfrac{Z^2}{\chi^2_\nu/\nu}$. Como $Z\sim N(0,1)$, $Z^2\sim\chi^2_1$ por definición de la chi-cuadrada. Por lo tanto,
	\begin{align*}
		T^2 = \frac{\chi^2_1/1}{\chi^2_\nu/\nu},
	\end{align*}
	que es exactamente la definición de $F_{1,\nu}$ (cociente de dos chi-cuadradas independientes, cada una dividida entre sus grados de libertad).
\end{solproblema}

% Evaluar
\begin{solproblema}[prob:e04a795]
	$F = 50/48 \approx1.042$. Como $1.042 < 3.36 = F_{14,11,0.025}$, \textbf{no se rechaza} $H_0:\sigma_1^2=\sigma_2^2$: las varianzas muestrales son consistentes con varianzas poblacionales iguales. En consecuencia, es apropiado usar la prueba $t$ \textbf{agrupada} (\emph{pooled}), que asume varianzas iguales y usa una varianza combinada $S_p^2$; no es necesario recurrir a la corrección de Welch, diseñada para el caso en que las varianzas son significativamente distintas — usar Welch aquí no sería incorrecto (es válida en general), pero sacrificaría innecesariamente algo de potencia estadística respecto a la prueba agrupada, que es más eficiente cuando su supuesto de igualdad de varianzas realmente se sostiene.
\end{solproblema}

% Crear
\begin{solproblema}[prob:6277cc2]
	\textbf{Ejemplo:} un laboratorio compara la precisión de dos instrumentos de medición de presión arterial. El instrumento A ($n_1=10$ mediciones repetidas) tiene $S_1^2=4.2\,\text{mmHg}^2$; el instrumento B ($n_2=8$) tiene $S_2^2=1.5\,\text{mmHg}^2$. El estadístico es $F = S_1^2/S_2^2 = 4.2/1.5 = 2.8$, con grados de libertad $(9,7)$. Comparando contra $F_{9,7,0.025}\approx4.82$: como $2.8 < 4.82$, no se rechaza $H_0:\sigma_1^2=\sigma_2^2$ al nivel $\alpha=0.05$; no hay evidencia suficiente de que la precisión de los dos instrumentos difiera.
\end{solproblema}
